\subsection{Vorgehen und Aufbau}\label{sec:vorgehen_und_aufbau}

Um das Konzept \enquote{Resteria} zu entwickeln, kombiniert die Projektgruppe zwei methodische Bausteine. Zum einen untersucht sie bestehende Plattformen aus dem Kulinarik- und Reste-Bereich in einer eigenen Wettbewerbsanalyse und leitet daraus die konkrete Marktlücke ab. Zum anderen stützt sie die Gestaltung einzelner Funktionen, etwa Gamification und Community-Mechanismen, auf verhaltenswissenschaftliche Literatur zu Nudges und nachhaltigem Verhalten.

Kapitel 2 bildet den Hauptteil dieses Berichts. Es beginnt mit der Wettbewerbsanalyse bestehender Plattformen, beschreibt anschließend die Zielgruppe genauer und entwickelt daraus das Konzept von Resteria mit seinen zentralen Funktionen wie Reste-Zutaten-Matching, Rettungspunkten und Reste-Level. Danach zeigt das Kapitel, wie Community-Aufbau und Feedback geplant sind, etwa über einen Rettungs-Feed und zeitlich begrenzte Challenges. Eine Phasenplanung mit Zeitrahmen und eine abschließende Reflexion zur Zielerreichung runden das Kapitel ab. Kapitel 3 fasst die Kernergebnisse zusammen, benennt einen Ausblick auf mögliche nächste Schritte und ordnet die Grenzen des Konzepts ein.
